% COURSE-ASSISTANT SOURCE — FALL 2025 ARCHIVE
% This material is not authoritative for Fall 2026; current-course material
% takes priority. See LN01_figure_notes.md for figures not stored here.
\documentclass[10pt]{beamer}

\input{EC2211_preamble.txt}
\title{Lecture Notes 1 \newline Measuring the Macroeconomy}


\begin{document}
\maketitle


\begin{frame}{Contents and literature}
\begin{itemize}
\item Measuring the macroeconomy
\item Mathematical preliminaries: logarithms and growth rates
\item Growth facts
\bigskip
\end{itemize} \medskip

Literature:
\small{
\begin{tightitemize}
    \item Jones (2024), chapters 1-3 and 6.5
\end{tightitemize}
}
\end{frame}


%%%%%%%%%%%%%%%%%%%%%%%%%%%%%%%%%%%%%%%%%%%%%%%%%%%
%%%%%%%%%%%%%%%%%%%%%%%%%%%%%%%%%%%%%%%%%%%%%%%%%%%
\section{Measuring the macroeconomy}

\begin{frame}{What is GDP and what does it measure?}
\begin{columns}
    \column{0.35\textwidth}
    \small{
    \begin{itemize}
         \item Consider this graph
        \item What do we need to know to make it "interesting"?
    \end{itemize}
    \bigskip
    }
    \column{0.65\textwidth}
    \begin{figure}
        \centering
        \includegraphics[width=0.94\linewidth]{Lect00a_Literature/GDP_se.png}
        \flushright{\tiny{Current prices. Source: Statistics Sweden.}}
    \end{figure}
\end{columns}
\end{frame}


\begin{frame}{GDP}
\begin{block}{GDP = Gross Domestic Product}
    The value of all final goods and services produced in an economy over a certain period
\end{block}
\bigskip
\small{
\begin{itemize}
    \item National accounts, developed since 1930's, international standards in the \href{https://unstats.un.org/unsd/nationalaccount/sna2025.asp}{United Nation's System of National Accounts}
    \item Three approaches: expenditure = income = production
\end{itemize}
}
\end{frame}


\begin{frame}{Expenditure approach}
\begin{block}{The National Income Identity}
    $Y = C + I + G + EX - IM$
\end{block}
\bigskip

\begin{table}
    \flushleft
    \begin{tabular}{rl}
         $Y$ & = GDP \\
         $C$ & = consumption \\
         $I$ & = investment \\
         $G$ & = government purchases \\
         $EX$ & = exports \\
         $IM$ & = imports \\
         \textcolor{gray}{$NX$} & \textcolor{gray}{= $EX - IM = $ net exports $=$ "trade balance"} \\
    \end{tabular}
\end{table} \medskip
\end{frame}


\begin{frame}{Example 1}
\begin{itemize}
    \item Homer and Marge grow 100 oranges \medskip
    \item ... sell to Edna, \$0.5 each \medskip
    \item If Edna eats all oranges:
    \begin{itemize}
        \item[]$Y = C = \$50$
    \end{itemize} \medskip
    \item If Edna eats 90 oranges and save 10 for next period:
    \begin{itemize}
        \item $C = \$45$ and $I = \$5$ (the increase in inventory is counted as an investment)
        \item Again $Y = C + I = \$50$
    \end{itemize} 
\end{itemize}
\end{frame}


\begin{frame}{Example 2}
\begin{itemize}
    \item Homer and Marge grow 100 oranges
    \item ... sell to Moe, \$0.3 each
    \item ... Moe sells to Edna, \$0.5 each
    \item Only \emph{final} expenditure counts, so $Y = C = \$50$
    \item (Alternatively, you can think of Moe as first making a positive investment when buying oranges, then disinvesting when selling the oranges)
\end{itemize}
\end{frame}


\begin{frame}{}
\begin{figure}
    \centering
    \includegraphics[width=0.9\linewidth]{LN01_Measuring_the_economy/GDP_SE_ExpApproach.png}
\end{figure}
\end{frame}


\begin{frame}{Expenditure approach to Swedish GDP in 2023}
\begin{columns}
\column{0.35\textwidth}
\small{
\begin{itemize}
    \item Consumption here refers to household consumption
    \item Investment here refers to private investment...
    \item ... and includes changes in inventories
\end{itemize}
\bigskip
}

\column{0.65\textwidth}
\footnotesize{
\begin{table}
    \centering
    \begin{tabular}{lrr}
    & SEK bn & \% of GDP \\
    \hline
     Gross domestic product  & 6,143 & 100 \\
     Consumption   & 2,798 & 46\\
     Investment & 1,241 & 20 \\
     Government purchases & 1,935 & 31 \\
     Net exports & 170 & 3 \\
     \emph{. exports} & \emph{3,391} & \emph{55} \\
     \emph{. imports} & \emph{3,220} & \emph{52} \\
     \hline
     \multicolumn{3}{l}{\footnotesize{Source: \href{https://www.statistikdatabasen.scb.se/pxweb/en/ssd/START__NR__NR0103__NR0103E/NR0103ENS2010T01NA/}{Statistics Sweden}}}
    \end{tabular}
\end{table}
}
\end{columns}
\end{frame}


\begin{frame}{Income approach}
\begin{block}{}
GDP = "wages" + "profits"
\end{block}
\bigskip

\begin{itemize}
    \item "Wages" = compensation to workers
    \item "Profits" = compensation to capital owners (interest payments, dividends)
\end{itemize}
\end{frame}


\begin{frame}{Example 1}
\begin{itemize}
    \item Homer and Marge grow 100 oranges \medskip
    \item ... sell to Edna, \$0.5 each \medskip
    \item GDP = Homer and Marge's income = \$50
\end{itemize}
\end{frame}


\begin{frame}{Example 2}
\begin{itemize}
    \item Homer and Marge grow 100 oranges \medskip
    \item ... sell to Moe, \$0.3 each \medskip
    \item ... Moe sells to Edna, \$0.5 each \medskip
    \item GDP = Homer and Marge's income + Moe's profit = \$30 + \$20 = \$50
\end{itemize}
\end{frame}


\begin{frame}{Income approach to Swedish GDP in 2023}
\footnotesize{
\begin{table}
    \centering
    \begin{tabular}{lrr}
    & SEK bn & \% of GDP \\
    \hline
     Gross domestic product  & 6,143 & 100 \\
     Compensation of employees & 2,938 & 48 \\
     \emph{. wages and salaries} & \emph{2,426} & \emph{39} \\
     \emph{. employers' social contributions} & \emph{511} & \emph{8} \\
     Operating surplus and mixed income & 1,980 & 32 \\
     Taxes on production and imports less subsidies & 1,225 & 20 \\
     \hline
     \multicolumn{3}{l}{\footnotesize{Source: \href{https://ec.europa.eu/eurostat/databrowser/explore/all/economy?lang=en&subtheme=na10.nama10.nama_10_ma&display=list&sort=category&extractionId=nama_10_gdp}{Eurostat}}}
    \end{tabular}
\end{table}
\bigskip
Note: The last row ("Taxes ...") consists mostly of VAT, import tariffs and payroll taxes. Income taxes and corporate profit taxes are \emph{not} in that item.
}
\end{frame}


\begin{frame}{Production approach}
\begin{block}{}
GDP = sum of value added in production
\end{block}
\bigskip

\begin{itemize}
    \item Only "new value" generated in each production step is counted \smallskip
    \item Consider example 2:
    \begin{itemize}
        \item Homer and Marge generate \$30 worth of value when they grow apples
        \item Moe generates an additional \$20 worth of value, maybe by advertising, packaging, bringing oranges closer to customers etc
    \end{itemize}
\end{itemize}
\end{frame}


\begin{frame}{Production approach to Swedish GDP in 2023}
\tiny{
\begin{table}
    \centering
    \begin{tabular}{lrr}
    & SEK bn & \% of GDP \\
    \hline
     Gross domestic product  & 6,143 & 100 \\
     Agriculture, forestry and fishing & 79 & 1 \\
     Mineral extract & 38 & 1 \\
     Manufacturing & 823 & 13 \\
     Electricity, gas, steam and air conditioning, water supply, waste & 173 & 3 \\
     Construction & 354 & 6 \\
     Wholesale and retail trade & 559 & 9 \\
     Transport and storage & 206 & 3 \\
     Hotels and restaurants & 93 & 2 \\
     Publishing, media, sound, telecom, programming, IT & 397 & 6 \\
     Financial services and insurance activities & 246 & 4 \\
     Real estate activities & 476 & 7 \\
     Legal, accounting and activities of head offices, management consultancy & 175 &  3 \\
     Architectural and engineering activities, R\&D & 138 &  2 \\
     Advertising and market research, veterinary activities & 58 & 1 \\
     Administrative and support service activities & 211 & 3 \\
     Education & 53 & 1 \\
     Human health, residential care, social work & 132 & 2 \\
     Arts, entertainment, recreation and other service activities & 86 & 1 \\
     Non-profit institutions servicing households & 67 & 1 \\
     Government & 1,123 & 18 \\
     Taxes on products less subsidies & 655 & 11 \\
     \hline
     \multicolumn{3}{l}{\tiny{Source: \href{https://www.statistikdatabasen.scb.se/pxweb/en/ssd/START__NR__NR0103__NR0103E/NR0103ENS2010T06NA/}{Statistics Sweden}}}
    \end{tabular}
\end{table}
\medskip
Note: The last row ("Taxes ...") consists mostly of VAT and import tariffs.
}
\end{frame}



%%%%%%%%%%%%%%%%%%%%%%%%%%%%%%%%%%%%%%%%%%%%%%%%%%%
%%%%%%%%%%%%%%%%%%%%%%%%%%%%%%%%%%%%%%%%%%%%%%%%%%%
\section{Measuring the macroeconomy: prices and quantities}

\begin{frame}{Do we produce more today or are prices just higher?}
\begin{figure}
    \centering
    \includegraphics[width=0.9\linewidth]{Lect00a_Literature/GDP_se.png}
    \flushright{\tiny{Current prices. Source: Statistics Sweden.}}
\end{figure}
\end{frame}


\begin{frame}{Prices and quantities: preliminaries}
\small{
\begin{itemize}
    \item Suppose that just two goods are produced: apples and oranges \pause
    \item Let $Q$ denote quantity and $P$ price
    \item In year 1, $Q^A_1=Q^O_1=100$ and $P^A_1=P^O_1=1$ \pause
    \item In year 2, $Q^A_1=100$, $Q^O_1=150$, $P^A_1=2$ and $P^O_1=1$ \medskip \pause
    \item \emph{Nominal} GDP is then:
    \begin{itemize} 
        \item[] $Y^N_1=1\times 100 + 1 \times 100 = 200 $ 
        \item[] $Y^N_2 = 2 \times 100 + 1 \times 150 = 350 $
    \end{itemize} \medskip \pause

    \item How much more do we produce in period 2? How much have prices increased? 
\end{itemize}
}
\end{frame}


\begin{frame}{Prices and quantities: preliminaries}
\small{
\begin{itemize}
    \item Use prices from period 1 to calculate "real" GDP:
    \begin{itemize}
        \item[] $Y^{P_1}_1 = 1 \times 100 + 1 \times 100 = 200 $
        \item[] $Y^{P_1}_2 = 1 \times 100 + 1 \times 150 = 250 $
    \end{itemize} \pause
    \item Use prices from period 2 to calculate "real" GDP in year 1:
    \begin{itemize}
        \item[] $Y^{P_2}_1 = 2 \times 100 + 1 \times 100 = 300 $
        \item[] $Y^{P_2}_2 = 2 \times 100 + 1 \times 150 = 350 $
    \end{itemize} \bigskip \pause
    \item The growth rate of real GDP is then
    \begin{itemize}
        \item[] $\frac{Y^{P_1}_2-Y^{P_1}_1}{Y^{P_1}_1} = 25\%$ if prices from period 1 are used \medskip
        \item[] $\frac{Y^{P_2}_2-Y^{P_2}_1}{Y^{P_2}_1} = 17\%$ if prices from period 2 are used
    \end{itemize}
\end{itemize}
}
\end{frame}


\begin{frame}{Prices and quantities: preliminaries}
\begin{itemize}
    \item Not obvious in this simple example how to separate quantities from prices \medskip
    \item These two methods are the \emph{Laspeyres} and \emph{Paasche} index methods, respectively \medskip
    \item In the real world, things are much more complicated \medskip
    \item Measures are note exact, methods vary between countries!
\end{itemize}
\end{frame}


\begin{frame}{How do we compare "today" with "yesterday"?}
\begin{columns}
    \column{0.5\textwidth}
    \begin{figure}
        \centering
        \includegraphics[width=0.9\linewidth]{LN01_Measuring_the_economy/1970s.png}
    \end{figure}

    \column{0.5\textwidth}
    \begin{figure}
        \centering
        \includegraphics[width=0.6\linewidth]{LN01_Measuring_the_economy/2020s.png}
    \end{figure}
\end{columns}
\smallskip
\end{frame}


\begin{frame}{Prices and quantities}
\small{
\begin{itemize}
    \item Difficult to separate price changes from quality improvements \medskip
    \item Consumption baskets change over time
    \begin{tightitemize}
        \item Our preferences change
        \item New goods and services are developed
    \end{tightitemize} \medskip
    \item Problems are more salient when comparison is over long time periods or if price movements are large \medskip
    \item A typical method used by statistical agencies is the \emph{chain weighted} index:
    \begin{tightitemize}
        \item "Average" of Laspeyres and Paasche
        \item Chain-linked over short time-intervals 
    \end{tightitemize}
\end{itemize}
}
\end{frame}


\begin{frame}{Nominal and real GDP: example terminology}
\small{
\begin{itemize}
    \item Nominal GDP: "current prices, million SEK" \bigskip
    \item Real GDP:
    \begin{tightitemize}
        \item "2015 prices, million SEK"
        \item "Constant prices, reference year 2024, SEK million"
        \item "Chain linked volumes, index 2010=100"
        \item "Chain linked volumes (2010), million euro"
    \end{tightitemize}
\end{itemize}
}
\end{frame}


\begin{frame}{Price indices: some examples}
\small{
\begin{itemize}
    \item GDP deflator: 
    \begin{tightitemize}
        \item GDP deflator $=$ nominal GDP $/$ real GDP
        \item Average price of goods produced in the economy, weighted by importance in production
    \end{tightitemize}
    \item CPI, "Consumer price index":
    \begin{tightitemize}
        \item Average price of goods consumed in the economy, weighted by importance in households' consumption baskets
    \end{tightitemize}
    \item The main difference between these relates to international trade
    \begin{tightitemize}
        \item Many goods in the consumption basket are imported
    \end{tightitemize}
\end{itemize}
}
\end{frame}


\begin{frame}{From 1970 to 2023}
\begin{itemize}
    \item[] GDP increased from SEK 196 bn in 1970 to SEK 6,143 bn in 2023: 
            \[ \frac{Y^N_{2023}}{Y^N_{1970}}=\frac{6,143}{196}\approx 31.4 \]
            \medskip \pause
    \item[] In addition to separating prices from quantities, we often also want to account for population growth:
    \[
    \frac{Y^N_{2023}}{Y^N_{1970}} = \frac{P_{2023}N_{2023}y_{2023}}{P_{1970}N_{1970}y_{1970}} = \frac{P_{2023}}{P_{1970}} \frac{N_{2023}}{N_{1970}} \frac{y_{2023}}{y_{1970}} 
    \]
    where $y$ is real GDP per person and $N$ the population size
\end{itemize}
\end{frame}


\begin{frame}{From 1970 to 2023}
\begin{figure}
    \centering
    \includegraphics[width=0.9\linewidth]{LN01_Measuring_the_economy/GDP_POP_se.png}
\end{figure}
\centering{\small{
    $\frac{P_{2023}}{P_{1970}} \approx 10.9$, 
    $\frac{N_{2023}}{N_{1970}} \approx 1.3$ and
    $\frac{y_{2023}}{y_{1970}} \approx 2.2$
}}
\end{frame}

\begin{frame}{From 1970 to 2023}
\begin{figure}
    \centering
    \includegraphics[width=0.9\linewidth]{LN01_Measuring_the_economy/GDP_POP_se_log.png}
\end{figure}
\centering{\small{This is the same graph, but now using a log scale for the y-axis}}
\end{frame}


\begin{frame}{Why logarithms?}
    \begin{columns}
        \column{0.5\linewidth}
        In a graph with a log scale for the vertical axis:
        \begin{itemize}
            \item Each vertical step represents the same \emph{relative} change
        \item A straight line means that the \emph{growth rate} is constant
        \end{itemize}

        \column{0.5\linewidth}
        \begin{figure}
            \centering
            \includegraphics[width=0.98\linewidth]{LN01_Measuring_the_economy/GDP_POP_se_log.png}
        \end{figure}
    \end{columns} \bigskip \pause

    The graph indicates that:
    \begin{itemize}
        \item The growth rates of prices and nominal GDP were high from 1970 to 1990
        \item There is no clear long-run trends in the growth rates of real GDP per capita or the population
    \end{itemize}
\end{frame}



%%%%%%%%%%%%%%%%%%%%%%%%%%%%%%%%%%%%%%%%%%%%%%%%%%%
%%%%%%%%%%%%%%%%%%%%%%%%%%%%%%%%%%%%%%%%%%%%%%%%%%%
\section{Mathematical preliminaries: logarithms and growth rates}

\begin{frame}{Mathematical preliminaries: the natural logarithm}

Economists use both {\color{red}$\ln x$} and {\color{red}$\log x$} to denote the natural logarithm of $x$. By definition:
\begin{equation*}
x=e^a \Leftrightarrow a=\ln x
\end{equation*}
\smallskip
\noindent Properties:
\begin{equation*}
\ln (xy) = \ln x + \ln y
\end{equation*}
\begin{equation*}
\ln (x/y) = \ln x - \ln y
\end{equation*}
\begin{equation*}
\ln x^\alpha = \alpha \ln x
\end{equation*}
and if $x$ is small:
\begin{equation*}
\ln (1+x) \approx x
\end{equation*}
\begin{equation*}
\ln (1+x)^\alpha \approx \alpha x    
\end{equation*}
\end{frame}


\begin{frame}{Growth rates}
Economists are often interested in growth rates. Suppose that $y$ grows at rate $g$, that is:
\begin{equation*}
    y_{t+1}=(1+g_{t+1})y_t
\end{equation*}\smallskip
Note from the properties of the logarithm that we then have:
\begin{equation*}
    \ln{y_{t+1}} = \ln{[(1+g_{t+1})y_t]} = \ln{(1+g_{t+1})} + \ln{y_t} \approx g_{t+1} + \ln{y_t}
\end{equation*}\smallskip
All this shows different ways that we can express the growth rate:
\begin{equation*}
    g_{t+1} = \frac{y_{t+1}-y_t}{y_t} \approx \ln{y_{t+1}}-\ln{y_t} = \ln{\left(\frac{y_{t+1}}{y_t}\right)}
\end{equation*}
\end{frame}


\begin{frame}{The constant growth rule}
Suppose that the growth rate $g$ is constant so that, for all $t$,
\[
y_{t+1}=(1+g)y_t
\]

Suppose that we start at $t=0$, with $y_0$, and that there are $3$ periods: \medskip
\begin{tabular}{r}
\\
  $y_1=(1+g)y_0$\\
  \\
  $y_2=(1+g)y_1$\\
  \\
  $y_3=(1+g)y_2$\\
\\
\end{tabular}
\end{frame}



\begin{frame}{The constant growth rule}
Substituting recursively:
\begin{align*}
y_3 &=(1+g)y_2\\ \\
	&=(1+g)\underbrace{(1+g)y_1}_{y_2}\\
	&=(1+g)(1+g)\underbrace{(1+g)y_0}_{y_1}\\
	&=(1+g)^3 y_0
\end{align*}

The constant growth rule:
\begin{equation*} 
y_t = (1+g)^t y_0
\end{equation*}
\end{frame}


\begin{frame}{Annualizing growth rates}
\begin{itemize}
    \item Suppose that we observe a variable $y$ in two years, $t$ and $s$

    \item Which constant annual growth $g$ rate takes $y$ from $y_t$ to $y_s$ in $s-t$ years?
        \[
            y_s = (1+g)^{s-t} y_t
        \]

    \item So the annualized growth rate is
        \[
            g = \left( \frac{y_s}{y_t} \right) ^\frac{1}{s-t} - 1
        \]
\end{itemize}
\end{frame}


\begin{frame}{Annualizing growth rates: example}
\begin{itemize}
    \item GDP per capita in Sweden increased from SEK 145,400 in 1950 to SEK 603,752 in 2024 {\footnotesize(constant prices, reference year 2024)} \medskip

    \item So:
        \[
            \frac{y_{2024}}{y_{1950}} \approx 4.15
        \] \medskip

    \item The average annual growth rate is
        \[
            g = \left( \frac{y_{2024}}{y_{1950}} \right) ^\frac{1}{2024-1950} - 1 \approx 0.0194 = 1.94 \text{ percent}
        \]
\end{itemize}
\end{frame}



%%%%%%%%%%%%%%%%%%%%%%%%%%%%%%%%%%%%%%%%%%%%%%%%%%%
%%%%%%%%%%%%%%%%%%%%%%%%%%%%%%%%%%%%%%%%%%%%%%%%%%%
\section{Measuring the macroeconomy: exchange rates and international comparisons}

\begin{frame}{How do we compare GDP across countries?}
\begin{itemize}
    \item GDP in 2019: 95.1 trillion yuan in China, \$20.6 trillion in the United States \smallskip
    \item Exchange rate 1 USD = 6.91 yuan \pause
    \begin{itemize}
        \item Convert Chinese GDP to dollars at market exchange rates: \smallskip
        \item[] \hspace{1cm} 95.1 trillion yuan $\times \frac{\$1}{6.91 \texttt{ yuan}} = \$13.8$ trillion 
    \end{itemize} \smallskip
    \item But the price level in China was around 68.4 percent of the price level in the United States
    \begin{itemize}
        \item Convert Chinese GDP to dollars using U.S. prices (i.e. at \mfemph{purchasing power parity}) \smallskip
        \item[] \hspace{1cm} 95.1 trillion yuan $\times \frac{\$1}{6.91 \texttt{ yuan} \times 0.684} = \$20.1$ trillion 
    \end{itemize} \smallskip
    \item So PPP adjusted GDP in China is roughly as high as that in the United States
\end{itemize}
\end{frame}



%%%%%%%%%%%%%%%%%%%%%%%%%%%%%%%%%%%%%%%%%%%%%%%%%%%
%%%%%%%%%%%%%%%%%%%%%%%%%%%%%%%%%%%%%%%%%%%%%%%%%%%
\section{Growth facts}


\begin{frame}
\frametitle{}
\bigskip
\begin{quote}
"It is a capital mistake to theorize before one has data. Insensibly one begins to twist facts to suit theories, instead of theories to suit facts."
\end{quote}
\bigskip
\begin{flushright}-Sherlock Holmes (A Scandal in Bohemia, 1891)\end{flushright}
\end{frame}

\begin{frame}
\frametitle{Growth facts}
\begin{itemize}
\item There is enormous variation in GDP per capita across countries \medskip
\item Growth rates vary substantially across countries \medskip
\item Growth rates also vary over time. Sustained growth is a recent phenomenon.
\end{itemize}
\end{frame}

\begin{frame}{Fact 1: There is enormous variation in GDP per capita across countries}
\small{
\begin{table}
    \centering
    \begin{tabular}{l r r}
         Country&GDP per capita &Percent of U.S.\\
         \hline
         United States&62,589&100\\
         Sweden& 52,433 &83.8 \\
         Germany&51,191 &81.8 \\
         Japan& 39,704 &63.4 \\
         Mexico& 18,737&29.9\\
         China& 14,129 &22.6 \\
         India&6,711&10.7\\
         Bangladesh& 4,658 &7.4 \\
         Chad& 1,615&2.6\\
         Burundi& 790&1.3\\
         \hline
         \multicolumn{3}{l}{\tiny{GDP per capita in 2019, in 2017 U.S. dollars. Source: PWT 10.01.}}
    \end{tabular}
\end{table}
}
\end{frame}

\begin{frame}{Fact 2: Growth rates vary substantially across countries}
\small{
\begin{table}
\resizebox{0.9\columnwidth}{!}{%
    \centering
    \begin{tabular}{l r r c r}
         &\multicolumn{2}{c}{Percent change in GDP per capita since}&&\multicolumn{1}{c}{Annual growth rate (\%)}\\ \cline{2-3} \cline{5-5}
         &1960&1990&& 1960-2019\\
         \hline
         United States&228&58&&2.0\\
         Sweden&289&73&&2.3\\
         Germany&359&82&&2.6\\
         Japan&625&35&&3.4\\
         Mexico&181&51&&1.8\\
         China&1,321&482&&4.6\\
         India&475&369&&3.0\\
         Bangladesh&198&204&&0.1\\
         Chad&12&35&&0.2\\
         Burundi&5&-22&&0.1\\
         \hline
         \multicolumn{5}{l}{\tiny{Last year of data is 2019. Calculations based on GDP per capita in 2017 U.S. dollars. Source: PWT 10.01.}}
    \end{tabular}
    }
\end{table}}
\end{frame}


\begin{frame}{Fact 3: Growth rates vary over time}
\small{
\begin{table}
\resizebox{!}{3.5cm}{%
    \centering
    \begin{tabular}{l r r}
         &\multicolumn{2}{c}{Annual growth rate (\%)} \\ \cline{2-3}
         &1960-1990&1990-2019\\
         \hline
         United States&2.5&1.6\\
         Sweden&2.7&1.9\\
         Germany&3.1&2.1\\
         Japan & {\color{blue}5.8} & {\color{red} 1.0}\\
         China&3.0&{\color{blue}6.3}\\
         Mexico&2.1&1.4\\
         India&{\color{red}0.7}&{\color{blue}5.5}\\
         Bangladesh&{\color{red}-0.1}&{\color{blue}3.9}\\
         Chad&-0.6&1.0\\
         Burundi&1.0&-0.9\\
         \hline
         \multicolumn{3}{l}{\tiny{Calculations based on GDP per capita in 2017 U.S. dollars. Source: PWT 10.01.}}
    \end{tabular}
    }
\end{table}}
\end{frame}


\begin{frame}
    \frametitle{Fact 3: Growth rates vary over time}
    Very little growth before 1800s
    \begin{itemize}
        \item Lack of good data
        \item Best estimates:
        \begin{tightitemize}
            \item[] ... no growth in first millennium A.D.
            \item[] ... then less than 0.1 percent per year to year 1800
        \end{tightitemize}
    \end{itemize}
    \bigskip

    Rapid growth thereafter
    \begin{itemize}
        \item U.S. GDP per capita has increased by factor 20 since year 1800
        \item ... but not everywhere
        \item ... and starting dates vary
    \end{itemize}
\end{frame}


\begin{frame}{Fact 3: Growth rates vary over time}
\begin{figure}
    \centering
    \includegraphics[width=0.9\linewidth]{Fig_LongRunGrowth.png}
\end{figure}
\tiny{Real GDP per capita, 2017 USD 1000s. Source: \href{https://www.rug.nl/ggdc/historicaldevelopment/maddison/}{Maddison project}}
\end{frame}


\begin{frame}{Rapid growth (both population and productivity) in "the West" from 1800s}
\begin{figure}
    \centering
    \includegraphics[width=0.95\linewidth]{WorldOutputShare.png}
\end{figure}
\footnotesize{Different regions' share of world output. Source: \href{https://www.core-econ.org/the-economy/microeconomics/01-prosperity-inequality-11-british-colonization-india.html}{COREecon Microeconomics, figure 1.18}}
\end{frame}



\begin{frame}{Growth rates vary over time, but relatively stable just below 2 percent per year in the United States}
    \begin{figure}
        \centering
        \includegraphics[scale=0.38]{Fig_1900sGrowth.png}
    \end{figure}
    \tiny{Real GDP per capita, 2017 USD 1000s, and trend lines. Source: \href{https://www.rug.nl/ggdc/historicaldevelopment/maddison/}{Maddison project}}
\end{frame}


\begin{frame}{U.S. population, GDP, and GDP per capita}
    \begin{figure}
    \includegraphics[scale=0.4]{MACRO6_FIG03.09.jpg}
    \end{figure}
    \flushright \tiny{Figure 3.9 in Jones (2024)}
\end{frame}


\begin{frame}{The great divergence}
    \begin{itemize}
    \item Cross-country differences in GDP per capita were relatively modest until 1700s \medskip
    \item Emergence of large cross-country differences in living standards from 1700 onwards \medskip
    \item But less so in recent decades 
    \begin{tightitemize}
        \item[] In 1990, 58 percent of the world's population lived in low-income countries
        \item[] ... today, only 9 percent (source: Gapminder)
    \end{tightitemize}
    \end{itemize}
\end{frame}


\begin{frame}{The distribution of world population by GDP per capita}
    \begin{figure}
    \includegraphics[scale=0.4]{MACRO6_FIG03.08.jpg}
    \end{figure}
    \flushright \tiny{Figure 3.8 in Jones}
\end{frame}


\begin{frame}{The share of the world population living in extreme poverty has fallen substantially}
    \begin{figure}
        \includegraphics[scale=0.55]{ExtremePoverty.png}
    \end{figure}
    \tiny{Figure 1.3 in \href{https://www.core-econ.org/the-economy/microeconomics/01-prosperity-inequality-04-income-inequality.html}{COREecon Microeconomics (2024)}}
\end{frame}



\begin{frame}{GDP is not welfare, but ...}
    \begin{figure}
        \includegraphics[scale=0.3]{Fig_gdp_lifeexp.png}
    \end{figure}
    \tiny Life expectancy at birth (2023 Human Development Index) and GDP per capita (PWT 10.01). Ireland and Qatar were dropped from the sample. \newline
    \tiny A more entertaining and dynamic version of this graph: \href{https://www.youtube.com/watch?v=jbkSRLYSojo}{Hans Rosling: 200 years in four minutes}
\end{frame}


\begin{frame}{Growth since 1960}
\begin{itemize}
    \item Some major economies have maintained or improved their position relative to the US (e.g. France, Italy, Spain and Japan)\medskip

    \item In the poorer group, some have improved (Korea, India, China) while others have fallen further behind (Nigeria, Ethiopia)\medskip

    \item Great heterogeneity: growth miracles and growth disasters abundant
\end{itemize}
\end{frame}




\begin{frame}{Questions}
\begin{itemize}
    \item Why are some countries so rich and others so poor?\medskip

    \item Why do some countries grow rapidly and some stagnate?\medskip

    \item Can growth continue in the long run? What are the sources and limitations of growth?
    \begin{tightitemize}
        \item Population growth? 
        \item Physical capital?
        \item Education, research, ideas?
    \end{tightitemize}
\end{itemize}
\bigskip
\centering{Malthus $\longrightarrow$ Solow $\longrightarrow$ Romer}
\end{frame}


\begin{frame}{What we did}
\begin{itemize}
    \item Measuring the economy: national accounts; real and nominal GDP; price indices
    \item Mathematical preliminaries: logarithms and growth rates
    \item Growth facts
\bigskip 
\item[] \textbf{Next time:} the production function; growth accounting, the Solow model
\end{itemize}
\end{frame}

\end{document}
